\subsection{Overview}

Subsurface (tile) drainage is a dominant conduit for nitrate-nitrogen (NO$_3$-N) transport from agricultural land, and the way a drainage system is designed and operated exerts strong, quantifiable control over how much nitrate reaches shallow groundwater, drain water, and downstream surface waters. The evidence assembled here spans continuously monitored field plots, drainage-district and watershed outlets, isotopic and gaseous tracer studies, and process-based simulation across North America, Europe, Asia, Australasia, and irrigated arid systems. Several conclusions are consistently supported. First, nitrate is the overwhelmingly dominant form of nitrogen exported in tile discharge, and tile flow itself typically accounts for a majority of watershed water and NO$_3$-N export \parencite{Williams2015,Grose2022,Hayashi1999}. Second, nitrate loads are governed first and foremost by hydrology: annual and event loads scale tightly with drainage flow, so that the great majority of export occurs during high-flow periods of the non-growing season \parencite{Kladivko2021,ROSSI1991,Bauwe2020,Ikenberry2014,Speir2021}. Third, drainage design (drain spacing and depth) and drainage \emph{management} (controlled drainage, subirrigation, and shallow drainage) modulate both flow volume and nitrate delivery, with managed strategies consistently reducing tile NO$_3$-N losses relative to conventional free drainage, primarily by reducing drainage volume \parencite{Cooke2012,Nash2015d,Elmi2004,Smith2019,Ross2016}. Fourth, edge-of-field practices (riparian buffers, integrated buffer zones, denitrification beds) and in-field management (cover crops, fertilizer rate and timing) provide additional, well-quantified reductions \parencite{Jaynes2014,Carstensen2021,Kaspar2007,Chaplot2004}. This section synthesizes the magnitudes, mechanisms, and management responses that the literature has established, together with the greenhouse-gas trade-offs and the modeling tools used to predict them.

\subsection{Nitrate Concentrations and Loads Under Conventional Free Drainage}

\subsubsection{Concentration magnitudes and drinking-water exceedances}

Tile-drainage water in intensively cropped landscapes carries persistently elevated nitrate. In three Iowa drainage districts, tile water averaged about 13 mg\,L$^{-1}$ NO$_3$-N (range 2.3--21.0), exceeding the 10 mg\,L$^{-1}$ drinking-water standard roughly 85--90\% of the time, with annual loads above 33 kg\,ha$^{-1}$\,yr$^{-1}$ \parencite{Schilling2012b}; a companion drainage-district study reported outlet concentrations of 0--29.3 mg\,L$^{-1}$ averaging 11.1--14.7 mg\,L$^{-1}$ \parencite{Schilling2013}. Across three tile-drained Iowa catchments, the flow-weighted mean was 17.1 mg\,L$^{-1}$ with 74.2\% of samples exceeding the 10 mg\,L$^{-1}$ maximum contaminant level (MCL) \parencite{Ikenberry2014}. Early Iowa work similarly found flow-weighted NO$_3$-N of 21.0 ppm and individual samples of 2.3--44.2 ppm \parencite{BAKER1975}. Elsewhere, drain-water NO$_3$-N averaged about 10 mg\,N\,L$^{-1}$ (nitrate being $>$90\% of total N) in Japan \parencite{Hayashi1999}, 6.6 $\pm$ 4.4 mg\,L$^{-1}$ (2.3--17.3) in eastern Canada \parencite{Smith2011}, and 14.3 mg\,N\,L$^{-1}$ (volume-weighted) in a dryland Washington catchment \parencite{Kelley2017}. Very high concentrations characterize some European and irrigated systems: Italian drainage water averaged 128, 216, and 94 mg\,NO$_3$\,L$^{-1}$ across three years, mostly exceeding the 90 mg\,NO$_3$\,L$^{-1}$ surface-water limit \parencite{ROSSI1991}; Polish drainage water usually exceeded 70 mg\,L$^{-1}$ (maximum 98.5) \parencite{Dragon2016}; and Iranian drain effluent ranged from 0.5 to 28.6 mg\,L$^{-1}$ \parencite{DarziNaftchali2017}. In a German tile-drained catchment, 15\% of high-frequency data exceeded the drinking-water limit despite low soil-surface N surplus \parencite{Bauwe2020}. By contrast, a long-term Indiana experiment kept annual NO$_3$-N generally in the 4--9 mg\,L$^{-1}$ range, below 10 mg\,L$^{-1}$, through reduced fertilizer rates and cover cropping \parencite{Kladivko2021,Kladivko2004}. In organic (peat) soils, drainage NO$_3$-N was even higher, averaging 14.8--42.7 mg\,L$^{-1}$ \parencite{MILLER1979}.

\subsubsection{Load magnitudes and variability}

Annual nitrate loads are highly variable and strongly tied to water yield. Iowa tile systems exported roughly 15--40 kg\,N\,ha$^{-1}$\,yr$^{-1}$ \parencite{Kelley2017,Schilling2012b,Ikenberry2014}, with early measurements spanning 0--93 kg\,ha$^{-1}$ and losses "highly variable and strongly dependent on the amount of water lost" \parencite{BAKER1975}. An Ontario field exported 39.1 kg\,ha$^{-1}$ over the monitoring period (flow-weighted 7.0 mg\,L$^{-1}$) \parencite{Smith2011}; an Indiana site accumulated 30.7 kg\,ha$^{-1}$ \parencite{Williams2022}. Italian losses averaged about 214 kg\,NO$_3$\,ha$^{-1}$\,yr$^{-1}$ ($\approx$50 kg\,N\,ha$^{-1}$) \parencite{ROSSI1991}. In Japanese onion systems an estimated 179 kg\,N\,ha$^{-1}$ (58\% of applied N) was leached by subsurface drainage \parencite{Hayashi1999}, and Latvian sites lost on average 14.3\% of applied mineral fertilizer N (net 18.90 kg\,ha$^{-1}$), representing an economic loss of about 61 EUR\,ha$^{-1}$\,yr$^{-1}$ \parencite{Siksnane2017}. Under irrigated arid conditions, nitrate movement below the root zone reached 44--153 kg\,NO$_3$-N\,ha$^{-1}$\,yr$^{-1}$ depending on crop and fertilization \parencite{ROBBINS1980}. These magnitudes reflect that NO$_3$-N losses to subsurface drains arise chiefly from asynchrony between mineralization and crop uptake and from large pools of potentially mineralizable soil organic N \parencite{Cambardella1999}.

\subsection{Hydrologic Control of Nitrate Export}

\subsubsection{Flow--load coupling and high-flow dominance}

The single most robust finding across scales is that nitrate load is controlled by drainage discharge. Annual nitrate-N losses were strongly related to annual drainage flow, with differences among drain spacings well explained by flow ($p<.001$) \parencite{Kladivko2021}. Annual discharge had a strong positive relationship with annual NO$_3$-N losses at both plot and ditch scales in Germany, leading the authors to conclude that hydrology is the key control on nitrate export \parencite{Bauwe2020}. In Italy, NO$_3$ losses correlated strongly with drainage water ($r=0.93$) but only weakly with precipitation ($r=0.45$) \parencite{ROSSI1991}, and Canadian drain NO$_3$-N concentrations rose with flow rate ($R^2=0.94$) \parencite{Smith2011}. High-flow periods dominate export: more than half of total nitrate export in Iowa catchments occurred during the highest 10\% of daily discharges, and nearly all during the upper half of daily flows \parencite{Ikenberry2014}; in Indiana watersheds, 79--95\% of NO$_3$-N export occurred at the 60th--100th percentile of flows \parencite{Speir2021}; and in Iowa BMP catchments the highest 50\% of flows carried most of the load \parencite{Law2022}. High-frequency monitoring of a drained North Carolina field likewise showed the highest flows accounting for a disproportionate share of nitrate loss \parencite{Liu2020}.

\subsubsection{Seasonality}

Because loads track flow, they are strongly seasonal. About 70\% of annual drainage and nitrate-N load occurred in the November--April fallow season in Indiana \parencite{Kladivko2021,Kladivko2004}. Nitrate losses in Italy peaked in winter and early spring \parencite{ROSSI1991}, and in Iowa most NO$_3$-N losses occurred when crops were absent (November--May) \parencite{Cambardella1999}. Watershed and tile NO$_3$-N loading occurred mainly during larger winter and spring flow events \parencite{Williams2015}, and seasonal patterns extended to loads highest in late spring and fall \parencite{Chaplot2004}. Short high-drainage episodes can dominate: 21 kg\,N\,ha$^{-1}$ was lost in only 11 days during snowmelt in Japan \parencite{Hayashi1999}. Concentration seasonality is also common---lowest in summer and highest in fall in Canada \parencite{Smith2011}, and declining to $\leq$5 mg\,L$^{-1}$ in late summer and $<$2 mg\,L$^{-1}$ during extended dry periods in Iowa \parencite{Schilling2013}. In dryland Washington systems, sudden early-winter increases in tile discharge and nitrate followed about 150--160 mm of fall precipitation that saturated the profile and mobilized high-nitrate soil water \parencite{Kelley2017,Keller2008}. Summer concentration spikes during corn years had minimal effect on annual loads because drainage was low then \parencite{Kladivko2021}.

\subsubsection{Concentration--discharge and hysteresis behavior}

High-frequency records reveal characteristic concentration--discharge (C--Q) signatures. Annual NO$_3$-N export is largely chemostatic---yields increase proportionately with runoff \parencite{Speir2021}---and persistent chemostasis (log-log slope magnitude $<$0.04) has been documented in Illinois tile drainage \parencite{Hu2024}. Watershed NO$_3$-N concentration increased with discharge up to the 75th flow percentile and then reached an asymptote in Ohio \parencite{Williams2015}. Seasonally, chemostatic fall and spring behavior can shift to source or transport limitation in winter and summer \parencite{Speir2021}. Storm-event hysteresis is predominantly counterclockwise (anti-clockwise): 77\% of storms in Indiana \parencite{Speir2021}, most of 22 delineated events in North Carolina \parencite{Liu2020}, and $>$75\% of events (annual mean hysteresis index $-0.34$) in a long-term Ohio record \parencite{Liu2022b}, indicating dominance of distal nitrate sources during delivery. Flushing versus dilution was more evenly split, with the flushing index near zero overall and strong seasonality (dilution in winter/summer, flushing in spring/fall) \parencite{Liu2022b,Speir2021}. The Ohio analysis concluded that intensive subsurface drainage, rather than deep groundwater, was the dominant factor shaping nitrate C--Q relationships, and that the Richards--Baker flashiness index tracked the history of drainage installation \parencite{Liu2022b}. At the event scale, total storm nitrate load is controlled by the same factors that dictate tile runoff, which in turn emerges from a threshold in above- and below-tile antecedent storage and soil-water mobilization \parencite{Cain2022}. Storm nitrate yields were strongly driven by event runoff ($R^2$ up to 0.97), while antecedent precipitation was a poor predictor \parencite{Speir2021}. In some settings nitrate load, but not concentration, was related to hydrograph metrics \parencite{Williams2022}, and tile discharge showed threshold behavior once soil-moisture deficit exceeded $-17\pm3$ mm and water table rose above $0.84\pm0.03$ m \parencite{Williams2022}.

\subsubsection{Rapid tile response and preferential flow}

Tile nitrate responds within hours of water inputs. After irrigation, tile nitrate rose sharply within hours, stayed high for days, then dropped abruptly \parencite{Mohanty1998}, and modeling indicated preferential flow intercepted by the drain originated close to and above the tile \parencite{Mohanty1998}. Bypass (macropore) flow was identified as a major mechanism carrying most leached nitrogen in Japan \parencite{Hayashi1999}. In cracking clay soils, summer drought produced nitrate-rich salt crusts that were rapidly mobilized to tile drains by crack flow during the first post-drought rains \parencite{Colombani2016}. Pulses or waves of water with differing NO$_3$-N concentrations moving through the profile were invoked to explain abrupt concentration changes in tile water \parencite{BAKER1975}. Higher-drainage catchments routed more event water (peak event-water fractions of 50\% versus 12\%) and produced larger, earlier solute pulses than low-drainage catchments \parencite{Kennedy2012}.

\subsection{Drainage Design: Drain Spacing and Depth}

Drain spacing and depth are established design controls on both drainage volume and nitrate loss. More closely spaced drains produce greater drainage per unit area and greater nitrate-N loss than wider spacings \parencite{Kladivko2021,Kladivko2004}; in Indiana the 5-m spacing generally yielded the most drainage and the 20-m the least, and relative flow differences among spacings narrowed as the system aged \parencite{Kladivko2004,Kladivko2021}. Importantly, nitrate \emph{concentrations} generally did not vary consistently with spacing---loads differences were driven by flow \parencite{Kladivko2021,Kladivko2004}. Modeling in China showed nitrogen loss increasing with drain spacing (as a second-degree parabola) and linearly with drain depth, and increasing with wider/higher filters \parencite{Tao2021}. Shallower drains were more effective at removing excess water but caused more nitrate loss, and in an Iranian heavy soil spacings $>$40 m could not combat waterlogging; the D0.40L50 configuration was judged best from combined nitrate and water-loss viewpoints \parencite{Hashemi2020}. In salt-affected coastal clay soils, wider spacing enhanced the risk of nitrate leaching and groundwater pollution relative to narrower spacing, which removed nitrate faster and added less to groundwater \parencite{Singh2002,Singh2007,Singh2001c}. A meta-analysis found greater relative reductions in tile discharge under drainage water management with narrower spacing and deeper drains \parencite{Ross2016}.

\subsection{Managed Drainage: Controlled Drainage, Subirrigation, and Shallow Drainage}

Managed drainage strategies consistently reduce nitrate loss compared with conventional free drainage, and the literature establishes both the magnitude and the mechanism.

\subsubsection{Controlled drainage and drainage water management}

In Illinois, annual NO$_3$ loads from managed systems were always less than from free drainage, with flow reductions of 35--96\% and NO$_3$ load reductions of 37--79\% (averaging 61\% across sites) \parencite{Cooke2012}. A synthesis reported drainage water management reducing tile flow by about 46\% and tile NO$_3$-N loads by about 45--48\%, with the caveat that surface runoff can increase and that reductions arise mainly from reduced tile discharge rather than enhanced denitrification \parencite{Ross2016}. Reviews likewise cite controlled drainage reducing NO$_3$-N exports by about 48\% \parencite{Ward2018}. In Missouri claypan soils, managed drainage cut annual water drained by at least 73\% and reduced annual NO$_3$-N loss by 78--85\% in two of four years (e.g., 3.8 versus 25.5 kg\,N\,ha$^{-1}$, an 85\% reduction) \parencite{Nash2015d}. In China, controlled drainage significantly reduced groundwater NO$_3$-N, NH$_4$-N, and total phosphorus, with the highest NO$_3$-N reduction reaching 73.2\% \parencite{Cui2016}. Submerged pipe (raised-outlet) operation reduced NO$_3$-N losses by 24\% and 51\% at 20-cm and 40-cm submergence, respectively, relative to free outflow \parencite{Tao2021}.

\subsubsection{Subirrigation and water-table management}

Subirrigation (SI) lowered soil-profile NO$_3$-N by up to 50\% relative to free drainage and reduced drainage-water NO$_3$-N concentrations by 55--73\% \parencite{Elmi2004}. Seasonal mean drainage-effluent NO$_3$-N fell by 74\%, 55\%, and 64\% across three years, and losses dropped from 7.2 to 1.5 kg\,ha$^{-1}$ in one year \parencite{Elmi2004}. SI reduced residual soil nitrate by 19--45\% and produced more rapid nitrate decline with depth, with water-quality benefits accruing even when SI outflow occasionally exceeded free-drainage outflow---implicating enhanced denitrification and reduced loading rather than merely reduced volume \parencite{Elmi2002}. Water-table control that elevated the water table enhanced denitrification and N$_2$O evolution while reducing potential N transport to surface water, denitrification increasing as the water table became shallower \parencite{KLIEWER1995}. Controlled drainage--subirrigation combined with a cover crop reduced five-year tile NO$_3$ concentration and loss by about 47\% (from 9.45 to 4.99 mg\,N\,L$^{-1}$ and 102 to 53.6 kg\,N\,ha$^{-1}$) \parencite{Drury2014}.

\subsubsection{Shallow drainage}

Comparative trials show shallow drainage reduces flow and load strongly. In Canada, shallow drainage reduced flow by 38--64\% relative to deep/conventional drainage (controlled drainage reduced flow by 24--37\%), and controlled drainage reduced NO$_3$-N loading by 42.3--58.2\% in two of three years \parencite{Smith2019}. Shallow drainage had the lowest mean nitrate concentration and lowest loading among three systems and exceeded the drinking-water guideline least often \parencite{Smith2011b,Smith2019}, though shallow drainage sometimes reduced yields whereas controlled drainage trended toward higher yields \parencite{Smith2019,Smith2011b}. Locally repairing or blocking a drainage section measurably altered whole-system runoff and, when drains were blocked to simulate controlled drainage, reduced ammonia and bacteria concentrations \parencite{Wesström2015,RITTER1995}.

\subsection{Fertilizer, Cropping, and Cover-Crop Controls on Nitrate Loss}

Nitrogen input rate is a strong lever on nitrate load but has a weaker, less consistent influence on concentration than hydrology does. Modeling showed that changing N application by $-60$\% to $+40$\% did not affect flow or sediment but strongly affected NO$_3$-N loads, with 20/40/60\% input cuts reducing loads by 23/49/94\% \parencite{Chaplot2004}. Fertilizer rate above recommendations pushed weekly drainage NO$_3$-N above 10 mg\,L$^{-1}$ and increased annual losses to 4--64 kg\,ha$^{-1}$ in mineral soils \parencite{MILLER1979}. Yet several long-term records found that reducing fertilizer rate did not immediately reduce nitrate export---precipitation, not fertilizer rate, predicted annual flux at a Washington catchment \parencite{Kelley2017}, and NO$_3$-N temporal patterns were unrelated to fertilizer timing or amount in Iowa \parencite{Cambardella1999}, consistent with soil-N legacy effects that may delay management responses by years to decades \parencite{Bauwe2020}. Latvian data showed no correlation between fertilizer incorporated and N lost \parencite{Siksnane2017}. Incidental loss within 30 days of fertilizer application nonetheless accounted for 52\% of cumulative tile NO$_3$-N load at one Indiana site \parencite{Williams2022}, and fall soil NO$_3$-N was the strongest control on leachate NO$_3$-N ($r^2=0.76$) in North Dakota, where careful N management returned leachate and groundwater NO$_3$-N toward pre-irrigation levels \parencite{Derby2009}.

Nitrogen source and timing generally had modest tile-drainage effects: manure versus urea did not differ in NO$_3$-N losses over multiple years \parencite{Randall2000}, and a large ($\sim$1,000 site-year) analysis found no significant difference in tile N export by application timing or method \parencite{Ward2018}. Isotopic evidence, however, showed pulsed, fertilizer-driven shifts in the NO$_3$ source signature during the corn growing season \parencite{Hu2024}. Cropping system matters strongly: replacing corn--soybean with pasture reduced NO$_3$-N losses to $<$3.5 t\,yr$^{-1}$ \parencite{Chaplot2004}; potato rotations produced the highest losses, and a winter-wheat cover after potato cut flow-weighted nitrate by 30\% to below 10 mg\,L$^{-1}$ \parencite{RITTER1995}; and alfalfa in irrigated rotations contributed substantial residual nitrate to subsequent crops \parencite{ROBBINS1980}. Cover crops are among the best-quantified in-field practices: a rye winter cover crop reduced flow-weighted NO$_3$ concentrations by 59\% and loads by 61\% over four years, chiefly by lowering concentration rather than drainage volume, whereas a permanent gamagrass strip did not significantly reduce losses \parencite{Kaspar2007}. Cover crops reduced tile NO$_3$-N by 21--38\% (concentration) and 14--16\% (load) in Canada while modestly increasing drainage volume \parencite{Drury2014}, and reviews report cover-crop NO$_3$-N reductions of 13--94\% \parencite{Ward2018}. Long-term declines in Indiana NO$_3$-N were attributed to combined reductions in fertilizer rate, rotation and tillage changes, and a cover crop acting as a trap crop \parencite{Kladivko2004,Kladivko2021}.

\subsection{Sources, Transport Pathways, and Transformations}

\subsubsection{Isotopic source apportionment and denitrification}

Dual-isotope studies indicate that drainage nitrate is often derived largely from organic sources (manure and soil organic matter) via nitrification, with denitrification not always dominant. In Ontario, $\delta^{15}$N and $\delta^{18}$O relationships lacked the 2:1 denitrification slope, pointing to organic sources and continued mineralization/nitrification through all seasons \parencite{Smith2011,Smith2011b}. Controlled incubation and field isotope work in Illinois quantified oxygen-isotope exchange during nitrification (exchange fractions of 59\% fertilized versus 0\% non-fertilized) and showed tile-drainage NO$_3$ isotopes were highly responsive to discharge and fertilizer input, with strong $\delta^{15}$N--$\delta^{18}$O coupling in unfertilized periods \parencite{Hu2024}. Where denitrification is active, it is detectable isotopically and geochemically: along Polish groundwater flow paths nitrate declined while sulfate, hardness, excess N$_2$, $\delta^{15}$N (9.9--23.5 per mil), and $\delta^{18}$O (7.9--17.2 per mil) increased \parencite{Dragon2016}. Nested French watersheds showed $^{15}$N enrichment as nitrate declined, and nitrate budgets implied 34--82\% retention dominated by underground (not benthic) denitrification \parencite{Billy2011,Billy2013}. In Iowa drainageways, denitrification in slowly permeable soils could remove essentially all upland groundwater nitrate, providing a natural treatment system \parencite{Schilling2007}, and deep unoxidized loess had the highest denitrification potential, with a large lifetime subsoil NO$_3$ removal capacity \parencite{Cambardella1999}. In arid Idaho, comparison of root-zone leaching with drainage concentrations (mean 3.2 ppm) suggested considerable denitrification and/or dilution at the water-table interface \parencite{ROBBINS1980}.

\subsubsection{Preferential flow, hydrograph sources, and legacy}

Tile discharge is dominated by pre-event water. Three-component separation showed tile flow was 65\% groundwater and 29\% shallow soil water, with precipitation only 6\% \parencite{Williams2022}. At the watershed scale, shallow groundwater supplied 86\% of tile flow, and tile drains carried the highest NO$_3$-N of any source, contributing over 99\% of watershed-scale nitrate export during a monitoring budget \parencite{Grose2022}. Tile drainage acts as a rapid conduit that limits infiltration of solutes into the subsoil: for a 100 kg\,ha$^{-1}$ urea scenario, 39\% of nitrate-N was leached through tile drains rather than reaching the subsoil, and without drainage nitrate could reach the profile bottom and groundwater \parencite{Filipovic2015}. Repair of a Swedish drainage system increased the proportion of infiltrated NO$_3$-N delivered directly to the stream via the tile outlet rather than diffusely via groundwater \parencite{Wesström2015}. Landscape position strongly differentiates groundwater chemistry, with oxic, high-nitrate uplands (8.5--12 mg\,L$^{-1}$) contrasting with reduced, low-nitrate drainageways ($<$0.5 mg\,L$^{-1}$) \parencite{Schilling2007}. Nitrate concentrations were consistently higher in tile drainage water than in water below the root zone and decreased with depth below the drains \parencite{Randall2000}, while subsurface-drainage NO$_3$-N averaged only 20\% of shallow-groundwater concentrations and lagged by 1--2 years \parencite{Derby2009}. Because losses stem from asynchronous mineralization and uptake and large mineralizable-N pools, land-management effects may only become visible after years to decades owing to soil-N legacies \parencite{Cambardella1999,Bauwe2020}. Rapid nitrification of applied ammonium ($>$90\% converted to NO$_3$ within 72 h) provides a ready nitrate supply for leaching \parencite{Hu2024}.

\subsection{Watershed-Scale Integration and Downstream Fate}

Tile drainage is the primary pathway controlling watershed water and nitrogen export in drained landscapes. In Ohio, tile drainage accounted for 56\% of annual watershed discharge (37--74\%) and 44--82\% (mean 62\%) of watershed NO$_3$-N export, with strong month-to-month coupling between tile and watershed loads \parencite{Williams2015}. Tile drains contributed an average of 57\% of streamflow and dominated watershed NO$_3$-N export in Indiana, where three large county drains disproportionately delivered 69\% of watershed nitrate, obscuring field-scale cover-crop effects \parencite{Grose2022}. Nitrate delivery from intensively drained Iowa landscapes can be spatially and within-season homogeneous, with the largest drains contributing 70--87\% of exported water and nitrate \parencite{Coupe2020}. Scaling drainage-district losses to the receiving Boone River suggested drainage districts could account for essentially all river nitrate export in most years, underscoring drainage-district-scale management \parencite{Ikenberry2014}. Downstream, concentrations decline with drainage area---e.g., from tile outlets averaging $\sim$13 mg\,L$^{-1}$ to 8.4, 7.6, 6.4, and 4.2 mg\,L$^{-1}$ progressively downstream in Iowa \parencite{Schilling2012b}---and nitrate is generally lowest in forested sub-catchments and highest in drained croplands \parencite{Billy2013}. Substantial in-network retention (34--82\%) occurs via underground denitrification \parencite{Billy2011}, but in-stream processing immediately below tile outlets offered little potential for meaningful downstream reduction because its timing was poorly matched to seasonal loading \parencite{Schilling2012b}. Achieving downstream targets was estimated to require about a 55\% load reduction at the tile outlet \parencite{Schilling2012b}. Regional analyses show catchment-outlet nitrate flux driven by both discharge and concentration, whereas tile and open-drain fluxes are overwhelmingly discharge-controlled, and groundwater plus farm-scale drainage inputs explained about 60\% of downstream N-load variability \parencite{Goeller2019}. Long-term monitoring can detect improvement: Czech outlet N-NO$_3$ declined from 1,175 to 587 $\mu$g\,L$^{-1}$ between 2004 and 2018 \parencite{Fraindová2022}, and stacked catchment-scale BMPs reduced NO$_3$-N load by 60.3\%, most strongly at high flow \parencite{Law2022}. Paired-watershed designs at drainage-district tile outfalls are feasible, with minimum detectable change around 8\% for the best control--treatment pair and improved detection when restricted to the March--July drainage period \parencite{Schilling2013}.

\subsection{Edge-of-Field and Downstream Mitigation}

Beyond in-field practices, edge-of-field structures provide well-documented nitrate removal. Redirecting tile flow into a riparian buffer diverted 55\% of tile flow (over 18,000 m$^3$) as infiltration, raised the water table $\sim$35 cm, and removed nitrate at 1.01 $\pm$ 0.07 mg\,N\,L$^{-1}$\,m$^{-1}$; chloride showed no gradient, confirming denitrification rather than dilution drove the decline \parencite{Jaynes2014}. Integrated buffer zones reduced drainage total-N loss by 29\% and 71\% at two sites (removal of 77 and 237 g\,N\,m$^{-2}$\,yr$^{-1}$, mostly nitrate), with the filterbed accounting for 70--93\% of NO$_3$ removal; hydraulic short-circuiting via overflow pipes and preferential surface paths limited performance \parencite{Carstensen2021}. Denitrification (woodchip) beds followed Michaelis--Menten kinetics (half-saturation constant 7.2 mg\,N\,L$^{-1}$, maximum rate 7.1 mg\,N\,L$^{-1}$\,h$^{-1}$ at 23.5\,$^{\circ}$C, $Q_{10}=2.8$), with seasonal removal of 0.01--4.64 mg\,N\,L$^{-1}$\,h$^{-1}$ and very low greenhouse-gas surface emissions \parencite{Ghane2015}. Saturated riparian buffers removed nitrate without increasing N$_2$O emissions, keeping total N$_2$O similar to traditional buffers and below corn--soybean rotations, and reducing indirect downstream N$_2$O \parencite{Davis2019}. Reviews also credit vegetated ditches, gypsum amendment, blind/gravel inlets, and BMP combinations (a SWAT analysis suggested $\sim$60\% pollutant-load reductions) with nutrient mitigation, though effects on dissolved phosphorus can differ from nitrate \parencite{Ward2018}. Ephemeral drainageways can themselves function as natural nitrate treatment systems \parencite{Schilling2007}.

\subsection{Greenhouse-Gas Dimensions of Drainage and Nitrate}

Drainage nitrate is coupled to greenhouse-gas fluxes, both directly and indirectly. Drainage water is commonly supersaturated with dissolved N$_2$O, CH$_4$, and CO$_2$; in Japan, indirect emissions via subsurface drainage were 0.747 kg\,N$_2$O-N, 0.115 kg\,CH$_4$-C, and 132 kg\,CO$_2$-C\,ha$^{-1}$\,yr$^{-1}$---equivalent to 4.6\%, 58.3\%, and 2.5\% of direct emissions---with N$_2$O-N/NO$_3$-N ratios (0.00076--0.0105) below the IPCC factor of 0.015 \parencite{Sawamoto2003}. Dissolved N$_2$O emissions were largest in upland-rice plots where leached N was greatest (50--67\% of aboveground emissions), and EF5g values exceeded the IPCC default by more than threefold in some systems \parencite{Minamikawa2010}. The presence of a groundwater table increased total N$_2$O emissions via dissolved pathways even under zero fertilizer-N \parencite{Minamikawa2013}. Two-phase co-evolution of [NO$_3^-$], [N$_2$O], and [CO$_2$] during non-rainy periods traced heterotrophic denitrification, yielding drainage-weighted EF5g of 0.0084 \parencite{Minamikawa2015}. In Canada, dissolved N$_2$O losses in tile drainage were substantive---approaching surface N$_2$O emissions---strongly seasonal (up to 62\% in the non-growing season), and declined along tile length, implying an in-drain sink and likely underestimation of the IPCC drainage coefficient \parencite{Burton2024}. A global synthesis of 52 studies found subsurface drains had the highest mean emission factor among artificial waters, that EFs scale downward with increasing NO$_3$ loading, and that the stationary IPCC EF tends to overestimate N$_2$O production \parencite{Webb2021}.

Conversely, installing subsurface drainage sharply reduces methane in paddy systems by improving aeration: drained fields emitted 71\% less CH$_4$ than non-drained fields (3.5 versus 14.8 g\,CH$_4$\,m$^{-2}$), with faster post-drainage Eh recovery and lower pre-submergence soil moisture \parencite{Shiratori2007}. Drainage lowered methane and CO$_2$ production potential by $\sim$40\% through reduced substrate availability and more oxidative conditions, accelerating rice-straw decomposition \parencite{Furukawa2008}. Management and fertilizer choices modulate N$_2$O: polymer-coated urea reduced soil N$_2$O by 47\% and NH$_3$ volatilization by 73\% versus non-coated urea, whereas tile drainage itself did not affect gaseous N loss in a claypan system \parencite{Nash2015e}. Controlled drainage kept cumulative N$_2$O below 1\% of applied N across drainage types \parencite{Smith2019}, and integrated buffer zones acted as net sinks of waterborne N$_2$O while remaining net CH$_4$ sources \parencite{Carstensen2021}. In saline systems, cumulative N$_2$O rose with wider drain spacing and lower organic-fertilizer replacement (up to 17.6 kg\,ha$^{-1}$ under 0\% organic, 20-m spacing) \parencite{Zhai2023}. Adding tile drainage to diversion terraces and grassed waterways reduced CO$_2$ and N$_2$O emissions relative to non-tiled structures by regulating soil water \parencite{Driscoll2025}. Carbon export from drained fields is dominated by inorganic carbon through tiles ($\sim$96\% of total carbon export; $\sim$1107 kg\,ha$^{-1}$ dissolved carbon leached in one year), so reducing water throughput supports carbon retention \parencite{Schilling2023}.

\subsection{Saline/Sodic Reclamation and Specialized Cropping Systems}

In salt-affected soils, drainage design governs both reclamation and nitrogen dynamics. Sustained subsurface drainage lowered soil salinity from 35 to 4 dS\,m$^{-1}$ at 15-m spacing and improved rice yield (6.5 versus 1.8--1.9 Mg\,ha$^{-1}$ at wide spacings), with narrower spacing favoring nitrate removal and reduced groundwater loading, while wider spacing allowed ammonium losses and impaired yields \parencite{Singh2002,Singh2001c}. Total N loss ranged 3.7--23.5 kg\,ha$^{-1}$\,yr$^{-1}$ across spacings, and subsurface drainage caused roughly 3--20\% loss of applied N through effluent \parencite{Singh2002}. Discontinuing drainage caused nitrite to accumulate 18--34-fold over three years as salinity inhibited nitrification, whereas functioning drainage promoted microbial activity, nitrification, and ammonium storage on the clay complex \parencite{Singh2007}. In subsurface-drained paddies, adding a canola crop improved soil-N utilization and considerably decreased N losses relative to fallow, keeping health-risk quotients below thresholds \parencite{DarziNaftchali2017}. Combining subsurface drainage with organic--inorganic fertilization improved saline soils, crop yield, and quality: smaller spacing and higher organic fertilizer raised alfalfa and sunflower yields and reduced surface-soil salinity (e.g., 36--46\% reduction at close spacing), though higher organic fertilizer could increase surface salt content \parencite{Zhang2024d,Zhai2023}. Optimization identified 10-m spacing with 75\% organic plus 25\% inorganic N as best for sunflower under TOPSIS evaluation \parencite{Zhu2025}. In irrigated sugarcane, a field N balance attributed 43\% of N input to nitrate drainage loss, 13\% to ammonium drainage loss, 9\% to denitrification, and 31\% to uptake \parencite{Matinzadeh2017}, while sewage-irrigation modeling identified denitrification as the main loss pathway with no soil-profile N accumulation up to $\sim$66 mg\,L$^{-1}$ \parencite{Yang2008}.

\subsection{Modeling Nitrate Dynamics Under Drainage}

Process-based models reproduce drainage nitrate concentrations and loads well after site-specific calibration and are established tools for evaluating management scenarios. DRAINMOD-N II and RZWQM2 both satisfactorily predicted tile NO$_3$-N in Iowa, with DRAINMOD-N II performing slightly better (concentration NSE 0.87 versus 0.76; loss NSE 0.84 versus 0.64) and both yielding threshold N rates close to field observations \parencite{Du2017}. RZWQM2 independently tracked NO$_3$-N concentration across N rates and identified an MCL-meeting application rate consistent with long-term data \parencite{Qi2012}. DRAINMOD-N II reproduced daily nitrate variation in an Iranian heavy soil, though it overestimated total seasonal losses, and it ranked drainage designs by combined nitrate/water-loss criteria \parencite{Hashemi2020}. HYDRUS captured how rainfall intensity/timing, initial groundwater depth, soil-N content, and submerged-pipe height controlled N loss \parencite{Tao2021}. A field-scale system-dynamics model matched cumulative nitrate loss (168.71 versus 176.68 kg\,ha$^{-1}$) and identified mineralization and denitrification rate constants as the most sensitive parameters \parencite{Matinzadeh2017}. The parsimonious NIT-DRAIN model simulated hourly nitrate concentrations with high performance (e.g., $R^2=0.83$, NSE 0.83, KGE 0.86), but its accuracy depended heavily on discharge inputs, degrading sharply when simulated rather than observed discharge was used \parencite{Chelil2022}. A FILTER-type model reproduced water tables and drainage nitrate for sewage irrigation \parencite{Yang2008}, and a CHAIN\_2D formulation with bimodal hydraulic functions improved simulation of preferential tile-flow nitrate transport \parencite{Mohanty1998}. At larger scales, SWAT estimated flow, NO$_3$-N, and sediment with monthly mean errors under 10\% and quantified farming-practice effects \parencite{Chaplot2004}, and a mechanistic denitrification-bed model coupling Forchheimer, Michaelis--Menten, and Arrhenius terms predicted outflow nitrate satisfactorily ($R^2=0.86$, NSE 0.85) \parencite{Ghane2015}. Groundwater flow-and-transport modeling showed that drainageway denitrification at 0.015 day$^{-1}$ could reduce upland nitrate to below detection at the catchment outlet, and that adding tile drainage would increase catchment drainage by 22.5 mm\,yr$^{-1}$ and deliver an additional 12.3 kg\,yr$^{-1}$ of nitrate \parencite{Schilling2007}.